\subsection{二つの三価頂点を持つ七高辺型の排除と，再利用する有限化}
\label{ng51120:obstruction}

\[
 G_{5,11}=\{01,02,03,04,12,13,24\}
\]
は，正の整数局所 charts，marked clutching，非負 node，対向 $D$ の相殺の下で存在しない．
二価 cores は $3,4$，三価 cores は $1,2$，四価 core は $0$ である．
この排除には分母および scale の探索上限を用いない．
次の型へ再利用するため，有限化の仮定と無限族の保持を記録する．

\paragraph{向き別の reverse 上界．}
二価 core を $i$，他方の二価 core を $j$，
固有の高近傍を $h$，残る三価頂点を $k$ とする．
canonical order $(A,B,C,D)$ と ray $j$ の行を，列 $(0,h,k)$ の順で書くと
\[
\begin{array}{c|c|c}
\text{型}&(A,B,C,D)&\text{ray }j\text{ の行}\\\hline
\mathrm I&(0,h,j,k)&(r,-r-1,0)\\
\mathrm {II}&(h,0,j,k)&(-r-1,r,0)\\
\mathrm {III}&(0,h,k,j)&(-S,T,-L)\\
\mathrm {IV}&(h,0,k,j)&(T,-S,-L)
\end{array}
\]
となる．共通低 puncture の通常 reverse の絶対値は四以下である．
正例外 $\{1,q-1\}$ と負例外 $\{-1,-(q-1)\}$ を捨てずに比較する．
正の単位 entry $x=sS/r=1$ は，正枝では $S>r$，
負枝では pair equation の符号から一般に排除できる．
これにより I 型は $r\le3$，II 型は $r\le4$，IV 型は $S\le4$ となる．

正の III 型は先に除く必要がある．相手の node と三高 matching により，
両二価 cores が $(r,s,z,S,T)=(1,1,1,2,6)$ に固定される．
外側低辺の $D$ 相殺が $q_{01}+q_{12}=4$ を要求し，両分母三以上と矛盾する．
負の III 型は $w=S=T\le4$ であり，他の負枝も $w\le4$ となる．
$w^2/2<H(H+1)<2w^2$ から $H\le5,r\le9$ を得る．

正枝の相手別の議論は，上界を他の graph へ移す際に確認すべき箇所である．
I--II では相手 $(r',S',T')=(1,2,5)$，
I 側は $(r,S,T)=(3,4,6),(3,5,7)$ に限られる．
II--II では $r=r'=1$，$S,S'\le3$ となり，pair equation が $T,T'<14$ を与える．
II--IV は通常接着なら $T\le5$，負例外なら $L=3$，
$r<S<2r$ から $S\le7,T\le9$ となる．
IV の残る相手は II または負の III であり，それぞれ $T\le6,T\le5$ である．
従って全体は，数学的に証明された有限範囲
\[
 1\le r\le9,\qquad 2\le S\le7,\qquad S\le T\le13
\]
に含まれる．この範囲は全候補を含む上界であって，全ての点が接着可能という意味ではない．

\paragraph{全 scale を保持した点・半直線表示．}
$h=S-r,\delta=T-S$ とすると
\[
 Fsz+As+Bz+C=0,
\]
\[
 F=(2r+1)h-r^2(\delta-1),\quad
 A=h(r+1),\quad B=-r(h+\delta-1),\quad C=r(r+1)
\]
である．$F\ne0$ では
$(Fs+B)(Fz+A)=K=AB-FC$ の全 signed divisors を用いる．
$K=0$ の二直線と $F=0$ の一次方程式は，整数合同式を解いて全半直線として残す．
この graph では book と異なり $L\le4$ を仮定できず，$F=0$ は実際に生じる．
無限族を持つ reverse triples は
\[
\begin{array}{c|c|l}
(r,S,T)&\text{場合}&\text{全正整数 scale}\\\hline
(1,2,6)&F=0&s=2z-1,\ z\ge1\\
(2,6,12)&F=0&s=3a-2,\ z=2a-1,\ a\ge1\\
(2,3,6)&K=0&s=2a,\ z=1,\ a\ge1\\
(3,4,6)&K=0&s=3a,\ z=2,\ a\ge1
\end{array}
\]
である．後二つだけを共鳴として残す手続きでは不完全になる．
整除性と奇数分母条件による剰余類分割の後，一つの二価 core あたり
$90$ 点型と $78$ 全半直線型，計 $168$ 表示を得る．
向き I,II,III,IV の表示数はそれぞれ $57,61,8,42$ である．
これは幾何的模型の個数ではなく，無限族を含む完全なパラメータ表示数である．

\paragraph{最後の低辺と，独立監査の役割．}
低辺 $34$ の node を通る型対は $3592$ 個である．
外側低辺 $14,23$ で各二十一 node 分配を使い，三高の低 row の一正二負，
共通低 puncture の通常・例外接着を課すと $1395$ 配置が残る．
例外接着の要求分母と affine 分母の一致から，
定数不一致 $759$，parameter 非整数 $550$，計 $1309$ 配置を除く．
残る $86$ 配置では低辺 $34$ の残差 $D^3_4+D^4_3$ が
$72$ 配置で常に奇数，$14$ 配置で非零定数となる．
定数の内訳は $8$ が四配置，$28$ が一配置，$-4$ が九配置である．
自由 parameter の上限を設けず，全ての非負整数に対して相殺を排除する．

主検算 \nolinkurl{verify_g5_11_obstruction_v20.py} は
全 affine 行列，node 行，要求分母，残差を証明書に保存する．
独立した \nolinkurl{audit_g5_11_certificates_v20.py} は主生成器を import せず，
canonical row/pair 式と元の node 定義から全 $1395$ 配置を再構成する．
さらに整数性の primitive step を用いる別の生成法で
$90$ 点と $11$ primitive 半直線を得て，
主検算の $78$ 剰余類半直線がその全体を覆うことを有限周期ごとに照合する．
この監査は半直線の最初の有限点だけを照合するものではない．
無上界の結論は解析的 reverse 上界と全整数族の処理の組合せに基づく．
